\section{Thirteen integers: CRT elimination and a coprime spectator guard}
\label{sec:jones-seeded}

The polynomial specified below has exactly thirteen unrestricted integer
unknowns, exact total degree $11\,546\,112$, and a division-free circuit of
$520$ live binary operations ($328$ multiplications, $169$ additions and
$23$ subtractions). Its integer zero existence is equivalent to a finite
syntactic ZFC refutation. Thus, with the stated inherited compiler theorem,
ZFC proves that its lack of integer zeros is equivalent to
$\operatorname{Con}(\mathrm{ZFC})$. This is an ordinary mathematical proof,
not a new Lean certificate, a proof of consistency, or a global minimality
claim. Unlike the ten-integer companion, the arithmetic finalizer here uses
only the established Jones ratio interface, not Sun's 2026 bridge.

\subsection{Degree-separated kernels and the complete polynomial}
Use the original package-31 system of 22 quadratic residuals
$r_j(p)$ in 31 strictly positive coordinates. The exact data are
\texttt{split/parent\_source.json}; the coherent fixed port tuple is
$(\beta,K_0,M_C,M_F)=(\beta,K,M_C,M_{F,\mathrm{source}})$ from
\texttt{FIXED\_BINDING.json}. All four compressed little-endian payloads are
included and hash-authenticated. The inherited compiler theorem identifies
its positive common zeros with finite syntactic ZFC refutations. This is the
original tuple, not a mixture with a later packet's smaller coefficients.

Perform just two elementary integer row operations:
\[
 \varrho_{14}=r_{14}+r_2,\qquad
 \varrho_{15}=r_{15}+r_3,\qquad \varrho_j=r_j\quad(j\ne14,15).
\]
The inverses are subtraction of rows 2 and 3, so the common zero set is
unchanged. Write $\varrho_j=\sum_\alpha r_{j,\alpha}p^\alpha$ and
\[
 c_{j,\alpha}=\begin{cases}
 2r_{j,\alpha},&|\alpha|=0,1,\\
 (\prod_i\alpha_i!)r_{j,\alpha},&|\alpha|=2.
 \end{cases}
\]
Use the 31 positive weights $w_i$ and the 22 targets $q_j$ in
\texttt{split/split\_code.json}, and set
\[
 D_d(T)=\sum_{j,\,|\alpha|=d}c_{j,\alpha}T^{q_j-\alpha\cdot w}
 \quad(d=0,1,2).
\]
The supplied exact certificate establishes
\[
 W=\max_iw_i=116,\quad Q_*=\max_jq_j=2129,\quad
 (\deg D_0,\deg D_1,\deg D_2)=(1719,2027,1929).
\]
The leading coefficient of $D_2$ is exactly $1$; every kernel has internally
distinct exponents. There are 89 source records. Checking all 22,544
same-degree source/expansion interactions proves
\[
 [T^{q_j}]\bigl(D_2(T)U(T)^2+D_1(T)U(T)+D_0(T)\bigr)
 =2\varrho_j(p),\qquad U(T)=\sum_i p_iT^{w_i}.
\]
The degrees are kept in separate products. Thus a linear label may equal a
quadratic label without causing cross-degree contamination. No homogeneous
dummy coordinate or full Sidon condition is needed. Let
\[
 K_*=\max_d(\deg D_d+dW)=2161,\qquad E_*=K_*+W+2=2279.
\]
The fixed coefficient bound, checked from the actual original port bit
lengths, is $|c_{j,\alpha}|<A_*=2^{2134927}$.

The unknowns are exactly
\[
 (\epsilon,g,h,s,w,\phi,i,j,z_1,z_2,z_3,r,v)\in\mathbb Z^{13}.
\]
All following expressions are polynomial abbreviations. Fix
\[
 \kappa=2134940,\qquad a=2^\kappa=8192A_*,\qquad
 \eta=2^{17079510}=2^{8\kappa-10},
\]
and define
\begin{align*}
 b&=2\epsilon^2,& B&=ab^2,\\
 \ell&=\sum_iB^{w_i},& T_*&=\sum_jB^{q_j},\\
 U_c&=\ell+g,&\mathcal G&=D_2(B)U_c^2+D_1(B)U_c+D_0(B),\\
 \mathcal G_+&=\mathcal G+B^{2162},&N_0&=B^{117},\\
 L&=2bB^{2279},&N&=\eta bB^{1201},\\
 c&=bB^{116},&G_c&=32c,\\
 M_c&=N_0-1-(b-1)\ell,&J_c&=2\mathcal G_++T_*,\\
 S_c&=g+N_0J_c,&T_c&=M_c+N_0(B-2)T_*,\\
 R&=L^2+(S_c+T_c)L+(L-1-S_c).
\end{align*}
The extra monomial $B^{2162}$ lies above every target. It repairs positivity
on arbitrary signed $g$ without changing an extracted residual or adding a
new unknown. Every exponent in these formulas is fixed and nonnegative.

Use the Jones registers with the exact CRT reparameterization of $H$ below, and $\omega=w^2$:
\begin{align*}
 U&=N^2\omega,&Y&=N^2s,\\
 M&=RY,&P&=2M^2U,\\
 K&=R+1+h(P-1),&A&=M(U+1),\\
 B_1&=2R+1,&W_1&=b\omega,\\
 C&=B_1+W_1+\phi,&J&=4A-5,&
 Q&=3W_1C-2(W_1^2-1),\\
 D&=(A^2-1)C^2+1,&E&=iJDC^2,&F&=(A^2-1)E^2+1,\\
 G&=A+F(CD+1-A),&H&=C+F(B_1+jC),&I&=(G^2-1)H^2+1,\\
 a_0&=(P^2-1)K^2+1,&a_1&=DFI,\\
 u&=J,&t&=Q.
\end{align*}
Set
\[
 \delta=C-KY,\qquad \Delta=g(c-g)K^2-c^2\delta^2,
 \qquad (d_1,d_2,d_3)=(3s-1,3\phi+2,3\Delta-1).
\]
The new coprime-spectator stack is
\[
 C_0=a_0,\qquad B_j^*=1+4d_jC_{j-1}z_j^2,\qquad
 C_j=C_{j-1}B_j^*\quad(j=1,2),
\]
\[
 B_3^*=1+4d_3C_2a_1z_3^2,\qquad C_*=C_3=C_2B_3^*.
\]
Thus the last guard contains the second square candidate $a_1$ as a
spectator. This is the only change to the Pell sign stack.
Finally define
\begin{align*}
 Z&=\epsilon wi z_1z_2z_3a_0,&\Sigma&=Zu^5,&\rho&=t+ur,\\
 \Theta&=(\rho^3-u^2\rho-u^3)(\rho^2+23u^2),&
 X&=\Sigma v-\Theta,\\
 A_0&=\Sigma^2C_*,&A_1&=\Sigma^2a_1.
\end{align*}
The requested thirteen-integer equation is the following single polynomial:
\begin{equation}\label{eq:jones-seeded}
 \boxed{\mathcal F_{13}=(X^2-A_0-A_1)^2-4A_0A_1=0.}
\end{equation}
The literal division-free DAG is \texttt{generated/signed\_coprime\_split\_crt.json}.

\subsection{Signed-input positivity and sharp block bounds}
Assume only $\epsilon\ne0$, hence $b\ge2$ and $B=8192A_*b^2$.
Geometric-series estimates give
\[
 D_2(B)>\tfrac12B^{1929},\quad
 |D_1(B)|<2A_*B^{2027},\quad |D_0(B)|<2A_*B^{1719}.
\]
Completing half of the quadratic square gives, for every real $U_c$,
\begin{align*}
 \mathcal G_+
 &\ge\tfrac12 D_2(B)U_c^2+B^{2162}
      -\frac{D_1(B)^2}{2D_2(B)}-|D_0(B)|\\
 &>\tfrac14 B^{1929}U_c^2+B^{2162}
      -4A_*^2B^{2125}-2A_*B^{1719}\\
 &>\tfrac14 B^{1929}U_c^2+\tfrac12B^{2162}.
\end{align*}
The last line uses $4A_*^2/B^{37}<1/4$ and
$2A_*/B^{443}<1/4$. For every \emph{integer} $g$, $U_c=\ell+g$ is an
integer. Therefore
\[
 S_c>\tfrac12N_0B^{1929}U_c^2+U_c+N_0B^{2162}-\ell>0.
\]
Indeed $tU_c^2+U_c\ge0$ for an integer $U_c$ and $t\ge1$, while
$\ell<N_0$. This explicitly absorbs the signed linear term $g$, rather than
incorrectly deducing $S_c>0$ merely from $J_c>0$.

The mask satisfies $0<M_c<N_0$, so $T_c>0$ as well. Consequently
\[
 R=L^2+S_c(L-1)+T_cL+L-1>L^2.
\]
Also
\[
 N^2=2^{-20}a^{15}B^{2403}<B^{2418}<B^{4558}<L^2,
 \qquad b\le c\le N.
\]
Thus $R\ge N\ge8$, $b\le R,N$, and all these estimates hold before any
sign guard, carry condition, or source-coordinate decoding.

For the range needed by the interval guard, assume $0<g<c$.
Since $\ell<2B^{116}$, one has $U_c<2bB^{116}$. Then
\[
 |\mathcal G|<12A_*b^2B^{2161},\qquad
 2|\mathcal G|<B^{2162}/256,\qquad T_*<B^{2162}/256.
\]
Writing $P_c=B^{2162}$, this gives
$0<J_c<(2+1/128)P_c$. Since $g<N_0$ and $P_c>128$,
\[
 0<S_c<3B^{2279},\qquad 0<T_c<B^{2279},\qquad
 S_c+T_c<4B^{2279}\le L.
\]
For the middle bound, use $(B-2)T_*<B^{2130}$ and
$1+B^{2130}<B^{2162}$. These sharp estimates are why the off-target padding
does not require increasing the block exponent from 2279 to 2280.

\subsection{Sparse carries and parity-tuned square threshold}
Suppose now $b=2^k$, so $B=2^\nu$ with $\nu=\kappa+2k$.
The low mask gives
\[
 g\mathbin{\&}M_c=0\quad\Longleftrightarrow\quad
 g=\sum_i n_i B^{w_i},\qquad0\le n_i<b.
\]
Here the binary ``and'' is proof notation only. Put $p_i=n_i+1$, so
$U_c=\sum_i p_iB^{w_i}$. Every coefficient $h_e$ of
$D_2(T)U(T)^2+D_1(T)U(T)+D_0(T)+T^{2162}$ obeys
\[
 |h_e|<994A_*b^2<B/8.
\]
The bound includes the possible collisions at unmarked coefficients and the
extra coefficient 1; it follows from $\sum_i p_i\le31b$.
The unnormalized digits of $J_c$ are $2h_e+\varepsilon_e$, where
$\varepsilon_e=1$ exactly at the targets. Starting with incoming carry 0,
normalized carries remain in $\{-1,0\}$. At a target, its digit avoids
$B-2$ exactly when it is 0 or 1. This happens exactly when $h_e=0$,
independently of an incoming borrow. The degree-separated extraction theorem
then gives, using the inverse row operations,
\[
 S_c\mathbin{\&}T_c=0\quad\Longleftrightarrow\quad
 p=n+1\text{ is a positive common zero of the original system}.
\]
Conversely every such zero has this code with $0<g<c$ for arbitrarily large
$b=2^{1+2r}$. The all-one positive tuple is not a zero, since original row 5
has value 1 there. Thus the encoded nonnegative tuple is not zero.

Write $L=2^\lambda$. Since $S_c,T_c>0$ and $S_c+T_c<L$, the three
nonoverlapping blocks in $R$ give
\[
 s_2(R)=\lambda+s_2(T_c)+1-\tau_2(S_c,T_c),
\]
where $\tau_2(S,T)=s_2(S)+s_2(T)-s_2(S+T)\ge0$ is the number of binary
carries. The 31 low allowed positions and 22 high target positions give
\[
 s_2(T_c)=139\nu-31k-22,\qquad\lambda=1+k+2279\nu.
\]
Therefore
\begin{align*}
 \lambda+s_2(T_c)+1
 &=2418\kappa+4806k-20\\
 &=2\bigl((8\kappa-10)+k+1201(\kappa+2k)\bigr)=\log_2N^2.
\end{align*}
Legendre's identity $v_2\binom{2R}{R}=s_2(R)$ now proves
\[
 N^2\mid\binom{2R}{R}\quad\Longleftrightarrow\quad
 S_c\mathbin{\&}T_c=0.
\]
The factor $b$ in $N$ is necessary in this parity class. It is a polynomial
factor, not a new witness or a variable exponent.

\subsection{An exact polynomial CRT reparameterization}
The inherited ratio interface uses $H_{\rm old}=B_1+j_{\rm old}C$ and
$JF\mid QF+J(H_{\rm old}-C)$. Since $E=iJDC^2$,
\[
 F\equiv1\pmod J,\qquad F\equiv1\pmod C,
\]
and hence $\gcd(J,F)=1$. The divisibility is equivalent to
$J\mid Q$ and $F\mid H_{\rm old}-C$. Our changed register
\[
 H=C+F(B_1+jC)
\]
automatically satisfies the latter. It also has the original form, with
\[
 j_{\rm old}=1+(A^2-1)i^2J^2D^2C^3B_1+Fj,
\]
which is a polynomial integer, since $F-1=(A^2-1)i^2J^2D^2C^4$.
This identity holds for arbitrary signed inputs, including $C=0$.

Conversely, at the canonical positive ratio witnesses, $0<B_1<C$ and
$j_{\rm old}\ge1$. Therefore $z=(H_{\rm old}-C)/F>0$ and
$z\equiv B_1\pmod C$, so $z=B_1+jC$ with $j\ge0$. Thus every
canonical witness is represented. No extra unknown or cofinal restriction is
introduced. The merger now has $u=J$ rather than $JF$. This CRT step is
inherited from the separately identified sparse-Pell packet, not claimed as
an original theorem here.

\subsection{The interval guard and the Jones interface}
If $\Delta>0$ and $b\ge2$, then $K\ne0$ and
$0<g<c$. The identity $4g(c-g)=c^2-(c-2g)^2\le c^2$ implies
$4\delta^2<K^2$. Along with $s\ge1$, $\phi\ge0$, and $w\ne0$, this gives
$C,Y>0$ and $K>0$. The usual square-gap argument forces $h\ge1$:
if $h\le0$, then $1\le K\le R+1<P$; $K=1$ would force $P-1\mid R$,
which is impossible, while for $2\le K<P$,
\[
 (PK-1)^2<(P^2-1)K^2+1<(PK)^2.
\]
Thus $a_0$ cannot be a square in that branch.

For precision, instantiate the free ceiling in
\texttt{Jones1982.Nine.ratio\_sound} by
$G_{\mathrm{aux}}=8gK^2+1$, rather than assuming its old margin condition.
The positive integer $K^2-4\delta^2$ then gives
\[
 (K^2-4\delta^2)G_{\mathrm{aux}}-8gK^2>0.
\]
This ceiling is a proof-side abbreviation; the ratio registers do not
otherwise depend on it.

Replace $i$ by $|i|$, which changes no condition. To handle signed $j$,
$u\mid t$ now means $J\mid Q$, while $F\mid H-C$ holds identically. Also
$F\equiv1\pmod D$, $G\equiv1\pmod D$, and $G\equiv A\pmod F$ show
that the positive integers $D,F,I$ are pairwise coprime.
Hence $a_1=DFI$ square makes each a square. Put $d=G^2-1>0$ and
$x=\sqrt I$. The norm-one unit $x+H\sqrt d$ is positive even if $H<0$.
A positive power of $G+\sqrt d$ is the identity modulo $CF$; multiplying by
sufficiently large such powers gives a Pell solution $H'>B_1$ with
$H'\equiv H\pmod{CF}$. Then $j'=(H'-B_1)/C>0$, all square conditions
remain true, and $J\mid Q$ remains true. The literal positive Jones interface
therefore applies and gives the power-of-two base and central-binomial
condition required above.

For completeness, take any common source zero, choose
$b=2^{1+2r}$ large enough, and set $\epsilon=2^r$. The direct source code
has $32g<G_c$, so the literal ratio completeness theorem supplies its
canonical witnesses. In \texttt{Ratio.lean}, \texttt{exists\_B25core} sets
\[
 \omega=2^{2R+1-k},\qquad U=(N^2/b)2^{2R+1}.
\]
Since $k=1+2r$ is odd and $k<2R+1$, this $\omega$ is a positive square.
No claim that the canonical $s$ is a square is needed.

The same canonical construction has
$\rho_0=(U+1)^{2R}/U^R=Y+f$ with
$0\le f<2^{2R}/U=b/(2N^2)$ and $Y\ge1$.
The source's \texttt{ratio\_core} gives
\[
 |C/K-\rho_0|\le\frac{2(Y+f)}{Y(U+1)}<\frac4U,
 \qquad |C/K-Y|<\frac b{N^2}.
\]
The last inequality uses $R\ge8$. Since $c,b\le N$ and
$g(c-g)\ge1$, it follows that
\[
 \frac{c^2(C-KY)^2}{K^2}<\frac{c^2b^2}{N^4}\le1\le g(c-g),
\]
so the very same canonical witnesses, translated by the exact CRT formula, satisfy $\Delta>0$.
This is a cofinal square-base proof, not an unsupported restriction of an
existing slack variable to squares.

\subsection{Raw two-root recovery, including the zero spectator branch}
For every integer assignment, $u=J=4A-5\ne0$. The polynomial
\[
 \Pi(x)=(x^3-x-1)(x^2+23)
\]
has no rational root and has a root modulo every positive integer.
For an odd prime other than $23$, either the cubic has a simple root or it
is irreducible, whose Frobenius three-cycle makes its nonzero discriminant
$-23$ a square. The quadratic then has a simple root. At $23$ the cubic has
the simple root $3$; at $2$, $-23\equiv1\pmod8$ admits compatible square
roots. Hensel lifting and CRT finish the assertion. Thus
$\Theta=u^5\Pi(t/u+r)\ne0$ for every assignment, and
\[
 \exists r\quad Zu^5\mid\Theta
 \quad\Longleftrightarrow\quad Z\ne0\ \hbox{and}\ u\mid t.
\]
Indeed a root modulo $|Z|$ proves the reverse implication once $u\mid t$.
If $v_p(u)=\alpha>\beta=v_p(t)$, the two factors of $\Theta$ have
valuations $3\beta$ and $2\beta$, contrary to $u^5\mid\Theta$.
If $Z=0$, divisibility would require $\Theta=0$, already impossible.

At a zero of the equation, $\Sigma=0$ is impossible because then the
polynomial is $\Theta^4\ne0$. In particular
$\epsilon,w,i,z_1,z_2,z_3,a_0$ are nonzero. Every $d_j\equiv2\pmod3$
is nonzero, and every $B_j^*\equiv1\pmod4$ is nonzero. Thus $C_2,C_*$
are nonzero.

\paragraph{No coincident radicands.}
One has $C_*\ne a_1$ without any positivity assumption. If $a_1=0$ this
follows from $C_*\ne0$. Otherwise $B_3^*\ne1$, and it cannot be $-1$
because it is $1$ modulo $4$. But
\[
 \gcd(B_3^*,a_1)=1,
\]
while $C_*=a_1$ would make $B_3^*$ divide $a_1$. This would force
$B_3^*=\pm1$, a contradiction.

Put $q=X/\Sigma$. The raw norm equation implies
\[
 (q^2+C_*-a_1)^2=4C_*q^2.
\]
If $q=0$, the norm would give $C_*=a_1$, just excluded. Therefore
\[
 r_0=\frac{q^2+C_*-a_1}{2q},\qquad r_1=q-r_0
\]
are rational, with $r_0^2=C_*$ and $r_1^2=a_1$. Rational square roots of
integers are integers. Hence $C_*,a_1$ are integer squares,
$q\in\mathbb Z$, and $\Sigma\mid\Theta$. In particular $J\mid Q$.
No complex-root separation assumption is used.

\paragraph{Recover the first two guards before ruling out $a_1=0$.}
The factors $a_0,B_1^*,B_2^*,B_3^*$ are pairwise coprime, because each
new factor is $1$ modulo its entire preceding product. Their nonzero product
$C_*$ is a square. Thus each absolute factor is a square. Since
$B_j^*\equiv1\pmod4$, every $B_j^*$ is positive: a negative square cannot
be $1$ modulo $4$. The product sign then forces $a_0>0$, also a square.
All prefixes are now positive, so $d_1,d_2>0$, yielding $s\ge1$ and
$\phi\ge0$.

At this point the last factor might still equal $1$ if $a_1=0$; this branch
must not be silently discarded. The unconditional source estimate above
has already proved $R>L^2>0$ for \emph{every} signed $g$. Since
$\epsilon,w\ne0$ and $s\ge1$, we have $b\ge2$, $N\ge8$, $U,Y>0$,
$A\ge2$, and, using $\phi\ge0$, $C>0$. Then
\[
 D\ge A^2>0,\qquad F\ge1,\qquad
 CD+1-A>0,\qquad G\ge A>1,\qquad I\ge1.
\]
Consequently $a_1=DFI>0$. Now $B_3^*\ne1$, and its positivity with
$C_2a_1z_3^2>0$ forces $d_3>0$, hence $\Delta>0$. This closes the
zero-spectator branch without adding $a_1$ to the nonzero multiplier.
The interval and ratio arguments above therefore apply.

\paragraph{Completeness.}
Start from the canonical positive ratio witnesses, translated by the CRT
parameterization. They have $a_0,a_1$ positive squares, $J\mid Q$, and all
three strict sign conditions. For $j=1,2$, $d_jC_{j-1}$ is positive
nonsquare because $d_j\equiv2\pmod3$ and $C_{j-1}$ is a square.
A nontrivial Pell solution $x^2-d_jC_{j-1}y^2=1$ can be chosen with $y$
even, so $z_j=y/2\ne0$ makes $B_j^*=x^2$. For the last guard,
$d_3C_2a_1$ is likewise positive nonsquare, since $C_2a_1$ is a positive
square. The same construction gives $B_3^*$ square and $C_*$ square.
The intersective lemma chooses $r$ with $\Sigma\mid\Theta$. Finally choose
\[
 v=\Theta/\Sigma+\sqrt{C_*}+\sqrt{a_1}\in\mathbb Z.
\]
Then the raw norm vanishes. Both directions are proved. By the inherited
compiler theorem, integer-zero existence is equivalent to a finite ZFC
refutation, and ZFC proves its negation equivalent to
$\operatorname{Con}(\mathrm{ZFC})$.

\subsection{Exact degree, comparison, and verification scope}
The exact relevant register degrees are
\[
\begin{array}{c|rrrrrr}
 &a_0&a_1&C_2&C_*&\Sigma&X\\\hline
 \deg&261250&2635110&1045009&4875914&448571&448572
\end{array}
\]
and $\deg A_0=5773056$, $\deg A_1=3532252$. The degree of $A_0$
strictly exceeds those of $X^2$ and $A_1$, so the output has leading term
$A_0^2$. Its diagonal leading coefficient is
\[
 116361\,a^{2885464}\eta^{960}\pmod{1000003},
\]
which is nonzero since $a$ and $\eta$ are powers of two. The circuit upper
bound and this specialization prove exact degree $11\,546\,112$.

The previous seeded thirteen-integer construction has degree $57\,866\,184$
and $759$ operations. At its unchanged old source, the coprime spectator
alone gives degree $46\,680\,720$ and $758$ operations, and adding CRT gives
degree $28\,847\,344$ and $754$ operations. With degree-separated extraction
and the one-monomial signed-input repair, the canonical result is degree
$11\,546\,112$ and $520$ operations: $80.05\%$ less degree and $31.49\%$
fewer operations than the seeded predecessor. All variants have thirteen
unrestricted integer unknowns and the same original coherent coefficient
tuple. The old-source proof is retained in
\texttt{old\_source\_section.tex} for an explicit apples-to-apples comparison.
No twelve-variable result, variable-count record, global optimum, or reduction
of coefficient storage is claimed.

The metric counts every live binary addition, subtraction and
multiplication, including scalar multiplication and all fixed-power gates.
Reuse and fixed leaves are free. The canonical evaluator uses sparse Horner
forms, reuses known exponent sums, and factors the merger cubic as
$\rho(\rho^2-u^2)-u^3$. These are exact polynomial identities, not new
semantic assumptions.

The checker independently parses all 22 original residuals, proves the two
row inverses and all 22 target polynomial identities, verifies the complete
22,544 isolation interactions, and authenticates the full fixed tuple. It
also replays literal circuit registers over three finite fields, checks signed
and zero branches, the norm and CRT identities, nonvacuous Pell examples,
source positivity, sparse carries, and the parity calibration. These finite
implementation audits corroborate the universal proof; they are not a new
kernel-checked formal theorem or an unbounded consistency search.

Run \texttt{python build\_jones.py}, \texttt{python check\_jones.py}, and
\texttt{python -O check\_jones.py}. The complete division-free polynomial and
fixed binding are in the emitted DAG and \texttt{FIXED\_BINDING.json}.
Original source provenance is in \texttt{SOURCE\_PINS.json}; the CRT source
archive and theorem are identified in \texttt{CRT\_PROVENANCE.json}.
